% Zone „Das kennst du schon“ – aus bank/strahlensaetze/zone.jsonl (zusammenbau v0.1)
\einheitenkopf[zone][Kennst du schon]{Das kennst du schon}
\setcounter{aufgabe}{0}
%% TODO Grundvorstellungs-Aufgabe als erste Hauptnummer der Zone (2.8) fehlt in der Bank

%% TODO Titel: Ich-kann-Satz fehlt in der Bank (Platzhalter: Kettenname) – Nr. 1
%% TODO Anweisung: ein Satz über der ersten Teilaufgabe fehlt in der Bank – Nr. 1
\begin{aufgabe}{Dezimalzahlen und ganze Zahlen mal und geteilt durch eine Zahl (auch nicht durch Zehnerpotenzen: dreißig geteilt durch fünfzig), Kommaverschiebung bei zehn, hundert, tausend}
%% TODO Darstellung für schwach fehlt in der Bank (strahlensaetze-zone-f1-v1; Abschnitt „Für schwache Schüler“)
\swz{$3{,}4 \cdot 100$ – wie viel?}{}
%% TODO Darstellung für schwach fehlt in der Bank (strahlensaetze-zone-f1-v3; Abschnitt „Für schwache Schüler“)
\swz{$2{,}4 \cdot 25$ – wie viel?}{}
\end{aufgabe}

%% TODO Titel: Ich-kann-Satz fehlt in der Bank (Platzhalter: Kettenname) – Nr. 2
%% TODO Anweisung: ein Satz über der ersten Teilaufgabe fehlt in der Bank – Nr. 2
\begin{aufgabe}{Längeneinheiten umrechnen (m in cm und zurück, km in m, mm in cm), Ergebnis in einer sinnvollen Einheit angeben}
%% TODO Darstellung für schwach fehlt in der Bank (strahlensaetze-zone-f2-v1; Abschnitt „Für schwache Schüler“)
\swz{$3\,\text{m}$ – wie viele Zentimeter?}{}
%% TODO Darstellung für schwach fehlt in der Bank (strahlensaetze-zone-f2-v3; Abschnitt „Für schwache Schüler“)
\swz{$4{,}8\,\text{km}$ – wie viele Meter?}{}
\end{aufgabe}

%% TODO Titel: Ich-kann-Satz fehlt in der Bank (Platzhalter: Kettenname) – Nr. 3
%% TODO Anweisung: ein Satz über der ersten Teilaufgabe fehlt in der Bank – Nr. 3
\begin{aufgabe}{Vielfache und Teiler, „doppelt“, „halb“, „dreifach“, „ein Drittel“ als Rechnung}
%% TODO Darstellung für schwach fehlt in der Bank (strahlensaetze-zone-f3-v1; Abschnitt „Für schwache Schüler“)
\swz{Das Doppelte von $35$ – wie viel?}{}
%% TODO Darstellung für schwach fehlt in der Bank (strahlensaetze-zone-f3-v3; Abschnitt „Für schwache Schüler“)
\swz{Ein Drittel von $4{,}5\,\text{m}$ – wie viel?}{}
\end{aufgabe}

%% TODO Titel: Ich-kann-Satz fehlt in der Bank (Platzhalter: Kettenname) – Nr. 4
%% TODO Anweisung: ein Satz über der ersten Teilaufgabe fehlt in der Bank – Nr. 4
\begin{aufgabe}{Dreisatz und fester Faktor („pro Portion“; k = Preis je Stück)}
%% TODO Darstellung für schwach fehlt in der Bank (strahlensaetze-zone-f4-v1; Abschnitt „Für schwache Schüler“)
\swz{$3$ Brötchen kosten $1{,}50\,€$. Was kosten $6$ Brötchen?}{}
%% TODO Darstellung für schwach fehlt in der Bank (strahlensaetze-zone-f4-v3; Abschnitt „Für schwache Schüler“)
\swz{Für eine Portion Pudding braucht man $0{,}25\,\text{l}$ Milch. Wie viel Milch brauchen $8$ Portionen?}{}
\end{aufgabe}

%% TODO Titel: Ich-kann-Satz fehlt in der Bank (Platzhalter: Kettenname) – Nr. 5
%% TODO Anweisung: ein Satz über der ersten Teilaufgabe fehlt in der Bank – Nr. 5
\begin{aufgabe}{Strecken auf Millimeter messen und zeichnen, Rechteck mit dem Geodreieck, Kreis mit dem Zirkel (Radius als halber Durchmesser), Parallelen mit dem Geodreieck zeichnen und erkennen}
\swa{Zeichne eine Strecke von $5{,}3\,\text{cm}$ Länge.}{\begin{ksys}[xmin=0,xmax=9,ymin=0,ymax=6]\end{ksys}}
\swa{Zeichne ein Rechteck mit den Seiten $6{,}5\,\text{cm}$ und $3{,}5\,\text{cm}$.}{\begin{ksys}[xmin=0,xmax=9,ymin=0,ymax=6]\end{ksys}}
\end{aufgabe}

%% TODO Titel: Ich-kann-Satz fehlt in der Bank (Platzhalter: Kettenname) – Nr. 6
%% TODO Anweisung: ein Satz über der ersten Teilaufgabe fehlt in der Bank – Nr. 6
\begin{aufgabe}{Draufsicht, Netz und Schrägbild eines Körpers unterscheiden; Draufsicht eines Zylinders ist ein Kreis, eines Quaders ein Rechteck}
\begin{teile}
\teil Eine Dose steht auf dem Tisch. Welche Form siehst du von oben? \\ \kreuz{Kreis} \\ \kreuz{Rechteck} \\ \kreuz{Dreieck}
\end{teile}
\begin{teile}
\teil Eine Rolle Küchenpapier liegt auf dem Tisch. Welche Form siehst du von oben? \\ \kreuz{Kreis} \\ \kreuz{Rechteck} \\ \kreuz{Quadrat mit Kreis}
\end{teile}
\end{aufgabe}

%% TODO Titel: Ich-kann-Satz fehlt in der Bank (Platzhalter: Kettenname) – Nr. 7
%% TODO Anweisung: ein Satz über der ersten Teilaufgabe fehlt in der Bank – Nr. 7
\begin{aufgabe}{Verhältnis lesen und als Division schreiben („drei zu fünf“ als Bruch und als Quotient, auch als Dezimalzahl)}
%% TODO Darstellung für schwach fehlt in der Bank (strahlensaetze-zone-f7-v1; Abschnitt „Für schwache Schüler“)
\swz{„drei zu vier“ – als Bruch?}{}
%% TODO Darstellung für schwach fehlt in der Bank (strahlensaetze-zone-f7-v3; Abschnitt „Für schwache Schüler“)
\swz{$15\,\text{cm}$ zu $25\,\text{cm}$ – als gekürzter Bruch?}{}
\end{aufgabe}

%% TODO Titel: Ich-kann-Satz fehlt in der Bank (Platzhalter: Kettenname) – Nr. 8
%% TODO Anweisung: ein Satz über der ersten Teilaufgabe fehlt in der Bank – Nr. 8
\begin{aufgabe}{Gleichung mit einer Variablen und Verhältnisgleichung nach x umstellen (x geteilt durch die bekannte Strecke gleich dem Verhältnis der zwei anderen; auf beiden Seiten mal nehmen)}
%% TODO Darstellung für schwach fehlt in der Bank (strahlensaetze-zone-f8-v1; Abschnitt „Für schwache Schüler“)
\swz{$\text{$x : 4 = 3$ – wie groß ist $x$?}$}{}
%% TODO Darstellung für schwach fehlt in der Bank (strahlensaetze-zone-f8-v3; Abschnitt „Für schwache Schüler“)
\swz{$\text{$\frac{x}{9} = \frac{2}{3}$ – wie groß ist $x$?}$}{}
\end{aufgabe}

%% TODO Titel: Ich-kann-Satz fehlt in der Bank (Platzhalter: Kettenname) – Nr. 9
%% TODO Anweisung: ein Satz über der ersten Teilaufgabe fehlt in der Bank – Nr. 9
\begin{aufgabe}{Winkel messen und vergleichen, Winkelsumme im Dreieck (zwei Winkel gleich → der dritte auch)}
%% TODO Darstellung für schwach fehlt in der Bank (strahlensaetze-zone-f9-v1; Abschnitt „Für schwache Schüler“)
\swz{Ein Dreieck hat die Winkel $50^\circ$ und $60^\circ$. Wie groß ist der dritte?}{}
%% TODO Darstellung für schwach fehlt in der Bank (strahlensaetze-zone-f9-v3; Abschnitt „Für schwache Schüler“)
\swz{Ein Dreieck hat die Winkel $38{,}5^\circ$ und $81{,}5^\circ$. Wie groß ist der dritte?}{}
\end{aufgabe}

%% TODO Titel: Ich-kann-Satz fehlt in der Bank (Platzhalter: Kettenname) – Nr. 10
%% TODO Anweisung: ein Satz über der ersten Teilaufgabe fehlt in der Bank – Nr. 10
\begin{aufgabe}{Längeneinheiten umrechnen (m in cm und zurück, km in m, mm in cm), Ergebnis in einer sinnvollen Einheit angeben – Fehler finden}
\swz[3]{Jonas rechnet: $0{,}7\,\text{m} = 7\,\text{cm}$. Finde den Fehler und rechne richtig.}{}
\end{aufgabe}

%% TODO Titel: Ich-kann-Satz fehlt in der Bank (Platzhalter: Kettenname) – Nr. 11
%% TODO Anweisung: ein Satz über der ersten Teilaufgabe fehlt in der Bank – Nr. 11
\begin{aufgabe}{Längeneinheiten umrechnen (m in cm und zurück, km in m, mm in cm), Ergebnis in einer sinnvollen Einheit angeben – selbst rechnen}
%% TODO Darstellung für schwach fehlt in der Bank (strahlensaetze-zone-f2-v7; Abschnitt „Für schwache Schüler“)
\swz{$0{,}4\,\text{m}$ – wie viele Zentimeter?}{}
\end{aufgabe}
